\documentclass{report}
\usepackage{listings}

\begin{document}
\title{labwork 2}
\maketitle

\section{Display the GPUs info}
The first step is to display the informations of the GPUs. It allows us to check if our hardware is well detected and to see the ids of our GPUs.
\begin{lstlisting}
GPUs = cuda.detect()
print(GPUs)
\end{lstlisting}
Output:
\begin{verbatim}
Found 1 CUDA devices
id 0    NVIDIA GeForce RTX 3060 Laptop GPU                              [SUPPORTED]
                      Compute Capability: 8.6
                           PCI Device ID: 0
                              PCI Bus ID: 1
                                    UUID: GPU-9e3ad670-7a54-1231-8562-a36df8ef464c
                                Watchdog: Enabled
             FP32/FP64 Performance Ratio: 64
                            FP16 Support: Yes
                            BF16 Support: Yes
                             FP8 Support: No
Summary:
        1/1 devices are supported
True
\end{verbatim}

\section{Display the name of the GPU}
Here we just want to show the name of the GPU. We select our GPU using the id that we just displayed.
\begin{lstlisting}
gpu = cuda.select_device(0)
print(gpu.name)
\end{lstlisting}
Output:
\begin{verbatim}
NVIDIA GeForce RTX 3060 Laptop GPU
\end{verbatim}

\section{Display the memory information}
The memory information are part of the Context. We retrieve a Tuple of values corresponding to the free memory and the total memory. Here, we only display the total. 
\begin{lstlisting}
print("total memory :", cuda.current_context().get_memory_info()[1])
\end{lstlisting}
Output:
\begin{verbatim}
total memory : 6441926656
\end{verbatim}

\section{Display the Cores information}
We want to get the number of cores but this information is not directly available. However, we can get the number of streaming multiprocessors of the GPU and the compute capability. Then with the dictionary from stackoverflow we can get the number of cores per streaming multiprocessors.
\begin{lstlisting}
# the following code is taken from: https://stackoverflow.com/questions/63823395/how-can-i-get-the-number-of-cuda-cores-in-my-gpu-using-python-and-numba
cc_cores_per_SM_dict = {
    (2,0) : 32,
    (2,1) : 48,
    (3,0) : 192,
    (3,5) : 192,
    (3,7) : 192,
    (5,0) : 128,
    (5,2) : 128,
    (6,0) : 64,
    (6,1) : 128,
    (7,0) : 64,
    (7,5) : 64,
    (8,0) : 64,
    (8,6) : 128,
    (8,9) : 128,
    (9,0) : 128,
    (10,0) : 128,
    (12,0) : 128
    }

my_sms = getattr(gpu, 'MULTIPROCESSOR_COUNT')
my_cc = gpu.compute_capability
cores_per_sm = cc_cores_per_SM_dict.get(my_cc)
total_cores = cores_per_sm*my_sms
print("GPU compute capability: " , my_cc)
print("GPU total number of SMs: " , my_sms)
print("total cores: " , total_cores)
\end{lstlisting}
Output:
\begin{verbatim}
GPU compute capability:  ComputeCapability(major=8, minor=6)
GPU total number of SMs:  30
total cores:  3840
\end{verbatim}

\end{document}
